% ============================================================================
% notes-preamble.tex : shared preamble for a self-contained lecture-notes set.
%
% This is the LINCHPIN. Author it FIRST, compile a 2-page smoke test, render
% the smoke test to PNG, and EYEBALL it BEFORE fanning out the build. A broken
% preamble fanned out in parallel wastes every parallel artifact. The smoke-test
% command is in SKILL.md; this file is meant to be \input by each notes.tex.
%
% Both (all) lecture-note documents for a topic \input or copy this block, so
% they render identically. Generalized from the LATS/MCTS notes; rename the
% colors, macros, and examplebox title to the new topic.
%
% PUNCTUATION: ASCII only. A repo write hook REJECTS Unicode em/en dashes
% (U+2014, U+2013). In LaTeX prose NEVER type a triple hyphen run (it renders an
% em dash) or a prose double hyphen run (it renders an en dash); use a comma,
% colon, parentheses, or the word "to". The one exempt double hyphen is a TikZ
% line segment inside a tikzpicture (an edge between two named coordinates,
% written backslash-draw, the source node, space, two hyphens, space, the
% target node). After compiling, verify the RENDERED pdf is dash-free:
%     pdftotext notes.pdf - | grep -c $'EMDASH-OR-ENDASH'   # must print 0
% ============================================================================
\documentclass[11pt]{article}

% math
\usepackage{amsmath}
\usepackage{amssymb}
\usepackage{dsfont}    % \mathds{1} indicator. THE \mathbb{1} TRAP: amssymb has
                       % NO blackboard-bold digit 1, so \mathbb{1} renders a junk
                       % (logic-looking) glyph. Use \mathds{1} for an indicator.
\usepackage{amsthm}
\usepackage{mathtools}

% figures, boxes, color
\usepackage{xcolor}
\usepackage{tikz}
\usetikzlibrary{arrows.meta, positioning, calc}
\usepackage[most]{tcolorbox}

% layout, links
\usepackage[margin=1in]{geometry}
\usepackage{hyperref}

% absorb the few hard-to-break inline-math lines (avoids tiny overfull boxes)
\setlength{\emergencystretch}{2em}

% a calm, low-saturation color scheme (rename per topic if desired)
\definecolor{calmblue}{RGB}{38, 70, 110}     % headings, structure
\definecolor{calmteal}{RGB}{38, 110, 102}    % definitions
\definecolor{calmsand}{RGB}{120, 95, 40}     % remarks, sidebars
\definecolor{calmgray}{RGB}{90, 90, 90}      % muted text
\definecolor{calmbg}{RGB}{245, 247, 250}     % box backgrounds

\hypersetup{
  colorlinks=true,
  linkcolor=calmblue,
  citecolor=calmteal,
  urlcolor=calmblue
}

% theorem-like environments.
% theorems, lemmas, propositions share one (italic) style and one counter.
\theoremstyle{plain}
\newtheorem{theorem}{Theorem}[section]
\newtheorem{lemma}[theorem]{Lemma}
\newtheorem{proposition}[theorem]{Proposition}

% definitions, examples, remarks share the (upright) definition style.
\theoremstyle{definition}
\newtheorem{definition}[theorem]{Definition}
\newtheorem{example}[theorem]{Example}
\newtheorem{remark}[theorem]{Remark}

% a tcolorbox to visually set off the running example (rename the title)
\newtcolorbox{examplebox}[1][]{
  colback=calmbg,
  colframe=calmblue,
  fonttitle=\bfseries\color{white},
  coltitle=white,
  title={Running example: REPLACE-WITH-YOUR-EXAMPLE-NAME},
  sharp corners,
  boxrule=0.6pt,
  #1
}

% TikZ styles shared by every figure. (The one place an ASCII double hyphen is
% legal is a line segment between two named coordinates inside a tikzpicture.)
\tikzset{
  treenode/.style={
    draw=calmblue, rounded corners, align=center,
    inner sep=3pt, font=\small, fill=white, minimum width=15mm
  },
  rootnode/.style={treenode, draw=calmteal, very thick},
  termnode/.style={treenode, fill=calmbg},
  ghost/.style={treenode, draw=calmgray, densely dashed, text=calmgray},
  treeedge/.style={-{Stealth[length=2mm]}, calmgray, thick},
  ghostedge/.style={-{Stealth[length=2mm]}, calmgray, densely dashed}
}

% convenience macros: define each KEY FORMULA ONCE so every document and every
% cross-reference renders the identical symbol. (Topic-specific; replace these
% with your topic's notation. The point is a single source of truth for numbers
% and notation: type the formula here, cite it everywhere.)
% \newcommand{\Foo}{...}
% \newcommand{\KeyFormula}{ ... typed once here ... }
% ============================================================================
% end shared preamble
% ============================================================================
